% Part: computability
% Chapter: machines-computations
% Section: halting-problem

\documentclass[../../../include/open-logic-section]{subfiles}

\begin{document}

\olfileid{tur}{und}{hal}
\olsection{Masalah Mandheg}

\begin{explain}
Upamane kita wis netepake sawijining enumerasi
andharan mesin Turing. Mula saben mesin Turing
nduwé \emph{indeks}: panggonane ing enumerasi
andharan mesin Turing $M_1$, $M_2$, $M_3$, \dots{}.

Kita ngerti yen mesthi ana fungsi kang ora bisa
diitung dening mesin Turing: himpunan andharan
mesin Turing---lan kanthi mangkono himpunan mesin
Turing---iku !!{enumerable}, nanging himpunan
kabeh fungsi saka $\Nat$ menyang $\Nat$ ora.
Nanging kita uga bisa nemokake tuladha fungsi
tartamtu kang ora bisa diitung. Salah sijine
yaiku fungsi mandheg.
\end{explain}

\begin{defn}[Fungsi mandheg]
 \emph{Fungsi mandheg}~$h$ ditegesi minangka
\[
h(e,n) =
\begin{cases}
  \text{0} & \text{yen mesin~$M_e$ ora mandheg kanggo input $n$} \\
  \text{1} & \text{yen mesin~$M_e$ mandheg kanggo input $n$}
\end{cases}
\]
\end{defn}

\begin{defn}[Masalah mandheg]
\emph{Masalah Mandheg} yaiku masalah nemtokake
(kanggo sembarang $e$, $n$) apa mesin Turing~$M_e$
mandheg kanggo input kang dumadi saka~$n$ guratan.
\end{defn}

\begin{explain}
Kita bakal nuduhake yen~$h$ ora bisa diitung
dening mesin Turing kanthi nuduhake yen fungsi
kang gegayutan, yaiku~$s$, ora bisa diitung
dening mesin Turing. Bukti iki gumantung marang
kasunyatan yen apa wae kang bisa diitung dening
mesin Turing uga bisa diitung dening mesin Turing
tertib (\olref[mac][dis]{sec}), lan yen rong mesin
Turing bisa disambung dadi siji mesin
(\olref[mac][cmb]{sec}).
\end{explain}

\begin{defn} Fungsi~$s$ ditegesi minangka
\[
s(e) =
\begin{cases}
  \text{0} & \text{yen mesin~$M_e$ ora mandheg kanggo input $e$} \\
  \text{1} & \text{yen mesin~$M_e$ mandheg kanggo input $e$}
\end{cases}
\]
\end{defn}

\begin{lem}
Fungsi~$s$ ora bisa diitung dening mesin Turing.
\end{lem}

\begin{proof}
Kanggo nggayuh kontradiksi, upamane fungsi~$s$
bisa diitung dening mesin Turing. Yen mangkono
ana mesin Turing~$S$ kang ngitung~$s$.
% OLPL-153: Sumber nyamake mandheg ing kothak kapisan karo definisi mesin tertib; teorema sadurunge marengake milih mesin tertib, kang banjur nduwe sipat kothak kapisan.
Tanpa nyuda umume, kita bisa milih mesin
iki kang tertib. Mula nalika $S$ mandheg,
sirah maca kothak kapisan. Mesin iki bisa
``disambung'' karo mesin liyane~$J$ kang
mandheg yen diwiwiti nganggo input~$0$
(yaiku yen maca $\TMblank$ ing kahanan
wiwitan nalika sirah ana ing kothak sisih
tengen simbol pucuk pita), lan yen ora,
terus mlaku menyang tengen tanpa mandheg.
$S \concat J$, mesin kang diasilake kanthi
nyambung $S$ karo~$J$, iku mesin Turing;
mula mesin iki $M_e$ kanggo sawijining~$e$
(yaiku ana ing enumerasi). Miwitiana $M_e$
nganggo input~$e$~$\TMstroke$. Ana rong
kemungkinan: $M_e$ mandheg utawa ora mandheg.
\begin{enumerate}
\item Upamane $M_e$ mandheg kanggo input
  $e$~$\TMstroke$. Mula $s(e) = 1$.
  Dadi $S$, nalika diwiwiti nganggo~$e$,
  mandheg kanthi output siji $\TMstroke$
  ing pita. Banjur $J$ diwiwiti nalika ana
  $\TMstroke$ ing pita. Ing kahanan iki
  $J$ ora mandheg. Nanging $M_e$ iku mesin
  $S \concat J$, mula tumindake kudu padha
  karo $S$ kang diterusake dening $J$
  (yaiku ing kahanan iki terus mlaku
  menyang tengen tanpa mandheg). Mula
  $M_e$ ora mungkin mandheg kanggo input
  $e$~$\TMstroke$.

\item Saiki upamane $M_e$ ora mandheg
  kanggo input $e$~$\TMstroke$.
  Mula $s(e) = 0$, lan $S$, nalika
  diwiwiti nganggo input~$e$, mandheg
  kanthi pita kosong. $J$,~yen diwiwiti
  nganggo pita kosong, langsung mandheg.
  Maneh, $M_e$ nindakake apa kang
  ditindakake $S$ banjur~$J$, mula
  $M_e$ mesthi mandheg kanggo input
  $e$~$\TMstroke$.
\end{enumerate}
Ing saben kasus kita tekan kontradiksi karo
panganggep kita. Mula mesin Turing~$S$
ora mungkin ana: $s$~ora bisa diitung
dening mesin Turing.
\end{proof}

\begin{thm}[Masalah Mandheg Ora Bisa Dirampungake]
\ollabel{thm:halting-problem} Masalah mandheg ora bisa
dirampungake, yaiku fungsi~$h$ ora bisa diitung
dening mesin Turing.
\end{thm}

\begin{proof}
Upamane $h$ bisa diitung dening mesin Turing,
umpamane dening mesin Turing~$H$. Kita bisa
nggunakake $H$ kanggo mbangun mesin Turing
kang ngitung~$s$: dhisik, gawe salinan input
(dipisahake dening simbol~$\TMblank$). Banjur
balia menyang wiwitan lan lumakna~$H$.
Kita mesthi bisa gawe mesin kang nindakake
langkah kapisan (delengen
\cref{tur:mac:dis:prob:copier}), lan yen $H$
ana, kita bisa ``nyambung'' mesin iku
karo mesin panyalin mau kanggo ngasilake
mesin anyar kang nemtokake apa $M_e$
mandheg kanggo input~$e$, yaiku ngitung~$s$.
Nanging kita wis nuduhake yen mesin kaya
mangkono ora mungkin ana. Mula $h$~uga
ora bisa diitung dening mesin Turing.
\end{proof}

\begin{prob}
Masalah Mandheg Telu (3-Halt) yaiku masalah
menehi prosedur putusan kanggo nemtokake apa
mesin Turing kang dipilih sembarang mandheg
nalika diwenehi input telung~$\TMstroke$
ing pita kang liyane kosong. Buktekna
yen masalah 3-Halt ora bisa dirampungake.
\end{prob}

\begin{prob}
Tuduhna yen masalah mandheg bisa dirampungake
kanggo pasangan mesin Turing lan input $M_e$
lan~$n$ kang $e \neq n$, mula masalah iku
uga bisa dirampungake kanggo kasus kang
$e = n$.
\end{prob}

\begin{prob}
Kita wis mbuktekake yen masalah mandheg
ora bisa dirampungake yen inpute wilangan~$e$
kang nemtokake mesin Turing~$M_e$ liwat
enumerasi kabeh mesin Turing. Kepriye yen
andharan mesin Turing saka
\olref[tur][und][enu]{sec} oleh langsung
digunakake minangka input? Apa ana mesin
Turing kang mutusake masalah mandheg nanging
nampa andharan mesin Turing minangka input,
dudu indeks? Jlentrehna apa sebabe bisa
utawa ora.
\end{prob}

\begin{prob} Tuduhna yen fungsi \emph{parsial}~$s'$ kang ditegesi minangka
  \[
  s'(e) =
  \begin{cases}
    \text{1} & \text{yen mesin~$M_e$ mandheg kanggo input $e$}\\
    \text{ora ditegesi} & \text{yen mesin~$M_e$ ora mandheg kanggo input $e$}
  \end{cases}
  \]
  \emph{bisa} diitung dening mesin Turing.
\end{prob}

\end{document}
